% =====================================================================
\chapter{Evaluation}
\label{ch:evaluation}
% =====================================================================

This chapter evaluates temporal fidelity in three steps: validation of
the measurement chain itself
(\cref{sec:eval:instrumentation,sec:eval:validation}), the baseline
comparison of stock Renode against hardware
(\cref{sec:eval:sleep,sec:eval:timer,sec:eval:quantum,sec:eval:workload}),
and the effect of the Temporal Fidelity Layer
(\cref{sec:eval:tfl}). \Cref{sec:eval:discussion} discusses findings
and threats to validity.

\section{Experimental Setup}
\label{sec:eval:setup}

All experiments use the configuration of \cref{sec:problem:methodology}:
a NUCLEO-F103RB running the timing-probe firmware (Zephyr~3.6.0,
\SI{1000}{\hertz} tick), captured with a Saleae Logic~8 at
\SI{100}{\mega\sample\per\second}; and Renode~1.15.3 executing the identical ELF with
\texttt{PerformanceInMips~72} and the default
\SI{100}{\micro\second} quantum unless stated otherwise. TFL runs use
the configuration of \cref{lst:repl}. Tool versions are pinned in
\cref{app:reproducibility}. Statistical procedures follow
\cref{sec:problem:methodology}: Mann--Whitney~$U$ at $\alpha = 0.01$,
rank-biserial $r$ as effect size, Shapiro--Wilk for normality.

\section{Instrumentation Overhead}
\label{sec:eval:instrumentation}

The GPIO marker itself costs time; if that cost were large or noisy it
would contaminate every measurement. \Cref{tab:instrumentation}
quantifies it: a set/clear marker pulse occupies a mean of
\SI{261}{\nano\second} (18 CPU cycles: 14 in the Zephyr GPIO driver, 4
in the bus access), with $\sigma = \SI{10}{\nano\second}$. All measured
effects in this chapter are two to four orders of magnitude larger, so
instrumentation overhead is negligible for the conclusions drawn.

\begin{table}[tb]
  \centering
  \caption[GPIO instrumentation overhead]{GPIO instrumentation overhead
    on the physical target (10\,000 marker pulses).}
  \label{tab:instrumentation}
  \small
  \begin{tabular}{@{}lll@{}}
    \toprule
    Metric & Value & Notes \\
    \midrule
    Min.\ HIGH pulse width & \SI{257}{\nano\second} & GPIO driver + bus call \\
    Mean HIGH pulse width  & \SI{261}{\nano\second} & systematic offset \\
    Standard deviation     & \SI{10}{\nano\second}  & measurement noise floor \\
    Equivalent CPU cycles  & 18                     & 14 driver + 4 bus \\
    \bottomrule
  \end{tabular}
\end{table}

\section{Measurement-Chain Validation}
\label{sec:eval:validation}

To bound the logic analyser's own error, a hardware timer PWM output
with a nominal \SI{1.000000}{\milli\second} period was captured
(\cref{tab:validation}). The measured mean deviates by only
\SI{20}{\nano\second} (timebase offset $2\times10^{-5}$) with
$\sigma = \SI{12}{\nano\second}$---confirming that the capture chain
resolves all effects studied here with ample margin.

\begin{table}[tb]
  \centering
  \caption[Logic-analyser validation]{Validation of the capture chain
    against a crystal-derived \SI{1}{\milli\second} PWM reference.}
  \label{tab:validation}
  \small
  \begin{tabular}{@{}ll@{}}
    \toprule
    Metric & Value \\
    \midrule
    Configured period  & \SI{1.000000}{\milli\second} \\
    Measured mean      & \SI{1.000020}{\milli\second} \\
    Timebase offset    & \SI{20}{\nano\second} \\
    Standard deviation & \SI{12}{\nano\second} \\
    \bottomrule
  \end{tabular}
\end{table}

\section{RTOS Sleep Timing (T1--T3)}
\label{sec:eval:sleep}

\begin{table}[tb]
  \centering
  \caption[Sleep-based periodic toggling, hardware vs.\ baseline
    Renode]{Sleep-based periodic toggling (T1--T3): physical hardware
    versus baseline Renode. Means in \si{\milli\second}, standard
    deviations in \si{\micro\second}.}
  \label{tab:sleep}
  \small
  \begin{tabular}{@{}lS[table-format=3.3]S[table-format=3.3]S[table-format=3.1]S[table-format=3.3]S[table-format=2.1]S[table-format=2.2]@{}}
    \toprule
    Test & {Nominal} & {HW mean} & {HW $\sigma$} & {Renode mean} & {Renode $\sigma$} & {$J$} \\
    \midrule
    T1 & 2.000   & 2.991   & 410.6 & 1.999   & 81.3 & 0.20 \\
    T2 & 20.000  & 20.981  & 400.2 & 19.992  & 83.8 & 0.21 \\
    T3 & 200.000 & 201.014 & 398.7 & 200.011 & 69.4 & 0.17 \\
    \bottomrule
  \end{tabular}
\end{table}

\Cref{tab:sleep} compares the toggle-period distributions produced by
\ksleep{}-based loops. Two effects stand out.

\paragraph{Systematic mean offset (hardware only).} Physical means
exceed nominal by ${\approx}\,\SI{1}{\milli\second}$ regardless of
period: Zephyr rounds sleep durations up to the next
\SI{1}{\milli\second} tick, and the phase of the request relative to
the tick adds a uniformly distributed extra delay. Renode reproduces
the tick mechanism but---because thread wakeup and tick processing are
quantum-aligned and phase-locked---the \emph{distribution} of that
extra delay collapses.

\paragraph{Suppressed variance (emulator).} Hardware exhibits
$\sigma \approx \SI{400}{\micro\second}$ across all three periods,
while Renode shows only \SIrange{69}{84}{\micro\second}: jitter ratios
of $0.17$--$0.21$. The differences are highly significant
(Mann--Whitney $p < 10^{-8}$, $|r| > 0.99$ for all three tests).
\Cref{fig:perioddist} visualises T2: the hardware distribution is broad
and right-shifted by the tick offset, whereas baseline Renode
concentrates tightly around the nominal value.

\begin{figure}[tb]
  \centering
  \includegraphics[width=\textwidth]{fig61_period_dist.pdf}
  \caption[Period distribution for \ksleep{}-based toggling
    (T2)]{Period distribution for \ksleep{}(\SI{10}{\milli\second})
    toggling (T2). Hardware (blue) is right-shifted by tick rounding
    and broad; baseline Renode (red) clusters unrealistically tightly
    at the nominal period (dashed).}
  \label{fig:perioddist}
\end{figure}

\section{Timer, Scheduling, and ISR Timing (T6--T9)}
\label{sec:eval:timer}

\begin{table}[tb]
  \centering
  \caption[Consolidated baseline comparison, all statistical
    scenarios]{Consolidated baseline comparison (hardware vs.\ stock
    Renode). $J$ per \cref{eq:jitterratio}; $p$ from Mann--Whitney~$U$.
    $J$ deviates from $1$ in both directions.}
  \label{tab:consolidated}
  \small
  \begin{tabular}{@{}lS[table-format=3.3]S[table-format=3.3]S[table-format=3.1]S[table-format=3.3]S[table-format=2.1]S[table-format=2.2]S[table-format=1.4]@{}}
    \toprule
    Test & {Nominal (\si{\milli\second})} & {HW mean} & {HW $\sigma$ (\si{\micro\second})} & {Ren.\ mean} & {Ren.\ $\sigma$ (\si{\micro\second})} & {$J$} & {$p$} \\
    \midrule
    T1 & 2.000   & 2.991   & 410.6 & 1.999   & 81.3 & 0.20  & 0.0000 \\
    T2 & 20.000  & 20.981  & 400.2 & 19.992  & 83.8 & 0.21  & 0.0000 \\
    T3 & 200.000 & 201.014 & 398.7 & 200.011 & 69.4 & 0.17  & 0.0000 \\
    T6 & 20.000  & 20.000  & 1.3   & 20.001  & 29.4 & 23.36 & 0.8510 \\
    T7 & 10.000  & 10.001  & 8.7   & 10.004  & 57.5 & 6.58  & 0.1941 \\
    T8 & 10.000  & 10.000  & 22.2  & 10.006  & 90.1 & 4.07  & 0.1632 \\
    T9 & 2.000   & 2.000   & 0.9   & 2.001   & 14.6 & 16.28 & 0.0077 \\
    \bottomrule
  \end{tabular}
\end{table}

For primitives driven by the hardware timer rather than the scheduler
tick, the fidelity failure inverts (\cref{tab:consolidated}, rows
T6--T9). A \ktimer{} callback that hardware delivers with
$\sigma = \SI{1.3}{\micro\second}$ jitters with
$\sigma = \SI{29.4}{\micro\second}$ in the emulator ($J = 23.4$),
because callback delivery snaps to quantum boundaries
(\srcS{4}) and translation-block granularity (\srcS{2}). ISR latency
(T9) shows the same pattern ($J = 16.3$). Notably, the emulated
\emph{means} remain accurate to within a few microseconds---the
deviation is almost purely in the variance, which is exactly the
component that mean-based validation would never detect.
\Cref{fig:cdf} shows the resulting cumulative distributions for T9: the
hardware CDF is a near-step function, the emulated CDF a wide ramp
spanning the quantum.

\begin{figure}[tb]
  \centering
  \includegraphics[width=\textwidth]{fig62_cdf.pdf}
  \caption[Cumulative distribution of ISR-driven period
    (T9)]{Cumulative distribution of the ISR-driven \SI{1}{\milli\second}
    toggle period (T9, two-toggle interval \SI{2}{\milli\second}).
    Hardware rises almost vertically ($\sigma=\SI{0.9}{\micro\second}$);
    baseline Renode spreads over tens of microseconds.}
  \label{fig:cdf}
\end{figure}

\section{Quantum Sensitivity}
\label{sec:eval:quantum}

Since \srcS{4} deviations scale with the quantum, one may ask whether
simply shrinking the quantum restores fidelity.
\Cref{tab:quantum} and \cref{fig:quantum} quantify the trade-off for
T6: reducing the quantum from \SI{100}{\micro\second} to
\SI{1}{\micro\second} improves the jitter ratio from $22$ to $0.23$
(overshooting into over-determinism), but inflates simulation
wall-clock time by two orders of magnitude. The default quantum is a
deliberate performance choice; hardware-level fidelity through quantum
reduction alone costs roughly $100\times$ simulation slowdown and still
cannot reproduce hardware's \emph{distribution shape}---only shrink
the emulator's artificial noise. This motivates TFL's approach of
\emph{modelling} the missing effects instead.

\begin{table}[tb]
  \centering
  \caption[Quantum sensitivity (T6)]{Renode scheduling-quantum
    sensitivity for timer-callback jitter (T6) and simulation cost of a
    \SI{90}{\second} firmware run.}
  \label{tab:quantum}
  \small
  \begin{tabular}{@{}S[table-format=5.0]S[table-format=4.2]S[table-format=4.2]S[table-format=2.3]@{}}
    \toprule
    {Quantum (\si{\micro\second})} & {T6 $\sigma$ (\si{\micro\second})} & {Jitter ratio $J$} & {Wall-clock (\si{\second})} \\
    \midrule
    10000 & 2932.80 & 2269.69 & 0.005 \\
    1000  & 301.52  & 233.35  & 0.045 \\
    100   & 28.60   & 22.13   & 0.450 \\
    10    & 2.87    & 2.22    & 4.500 \\
    1     & 0.29    & 0.23    & 45.000 \\
    \bottomrule
  \end{tabular}
\end{table}

\begin{figure}[tb]
  \centering
  \includegraphics[width=\textwidth]{fig64_quantum_tradeoff.pdf}
  \caption[Quantum fidelity--performance trade-off]{The
    fidelity--performance trade-off of the Renode scheduling quantum
    (log--log). Each decade of quantum reduction buys one decade of
    jitter reduction at one decade of simulation slowdown.}
  \label{fig:quantum}
\end{figure}

\section{Workload Interference (T10)}
\label{sec:eval:workload}

\Cref{tab:workload} and \cref{fig:workload} examine how competing
computational load degrades a \SI{10}{\milli\second} timer callback.
On hardware, jitter grows from \SI{1.2}{\micro\second} (idle) to
\SI{24.5}{\micro\second} (extreme load)---a $20\times$ degradation
driven by interrupt masking and bus contention. Baseline Renode starts
an order of magnitude too noisy (\SI{28.9}{\micro\second} idle) and
grows only $1.5\times$, badly compressing the dynamic range of the
load effect. TFL (third column) tracks the hardware trend far more
closely, staying within a factor of $\le 7$ at idle and $1.1$ at
extreme load.

\begin{table}[tb]
  \centering
  \caption[Workload impact on timer jitter]{Workload impact on
    \SI{10}{\milli\second} timer-callback jitter ($\sigma$, in
    \si{\micro\second}).}
  \label{tab:workload}
  \small
  \begin{tabular}{@{}clS[table-format=2.1]S[table-format=2.1]S[table-format=2.1]@{}}
    \toprule
    Phase & Load & {Hardware} & {Baseline Renode} & {TFL} \\
    \midrule
    1 & None    & 1.2  & 28.9 & 8.5 \\
    2 & Light   & 3.1  & 32.1 & 10.2 \\
    3 & Heavy   & 8.7  & 38.5 & 14.8 \\
    4 & Extreme & 24.5 & 44.2 & 26.1 \\
    \bottomrule
  \end{tabular}
\end{table}

\begin{figure}[tb]
  \centering
  \includegraphics[width=\textwidth]{fig65_workload_impact.pdf}
  \caption[Workload impact on timer jitter]{Timer-callback jitter under
    increasing computational load (T10). Baseline Renode compresses the
    hardware's $20\times$ dynamic range into $1.5\times$; TFL restores
    the trend.}
  \label{fig:workload}
\end{figure}

\section{Effect of the Temporal Fidelity Layer}
\label{sec:eval:tfl}

\begin{table}[tb]
  \centering
  \caption[Jitter fidelity, baseline vs.\ TFL]{Jitter fidelity before
    and after TFL. $J$ = baseline Renode jitter ratio,
    $J_{\mathrm{TFL}}$ = TFL-enhanced jitter ratio (target $1.0$).
    Improvements in \textbf{bold}.}
  \label{tab:tfl}
  \small
  \begin{tabular}{@{}llS[table-format=2.1]S[table-format=2.1]S[table-format=4.1]S[table-format=2.2]S[table-format=2.2]@{}}
    \toprule
    Test & Scenario & {HW $\sigma$ (\si{\micro\second})} & {Ren.\ $\sigma$ (\si{\micro\second})} & {TFL $\sigma$ (\si{\micro\second})} & {$J$} & {$J_{\mathrm{TFL}}$} \\
    \midrule
    T1 & \ksleep{} \SI{1}{\milli\second}    & 410.6 & 81.3 & 39.7   & 0.20  & 0.10 \\
    T2 & \ksleep{} \SI{10}{\milli\second}   & 400.2 & 83.8 & 418.4  & 0.21  & {\bfseries 1.05} \\
    T3 & \ksleep{} \SI{100}{\milli\second}  & 398.7 & 69.4 & 3991.0 & 0.17  & 10.01 \\
    T6 & \ktimer{} \SI{10}{\milli\second}   & 1.3   & 29.4 & 8.5    & 23.36 & {\bfseries 6.77} \\
    T7 & Thread scheduling                  & 8.7   & 57.5 & 11.6   & 6.58  & {\bfseries 1.32} \\
    T8 & Multi-thread interference          & 22.2  & 90.1 & 24.6   & 4.07  & {\bfseries 1.11} \\
    T9 & ISR latency                        & 0.9   & 14.6 & 1.5    & 16.28 & {\bfseries 1.63} \\
    \bottomrule
  \end{tabular}
\end{table}

\Cref{tab:tfl} presents the central result. With TFL enabled
(full configuration, \cref{lst:repl}), five of seven scenarios move
substantially toward $J = 1$:

\begin{itemize}
  \item \textbf{Timer/scheduler/ISR scenarios (T6--T9)}: the emulator's
        artificial quantum noise is replaced by calibrated Gaussian
        jitter, bringing $J$ from $4.1$--$23.4$ down to $1.1$--$6.8$;
        T7--T9 land within a factor of $1.7$ of hardware variance.
  \item \textbf{Sleep at \SI{10}{\milli\second} (T2)}: $J$ improves
        from $0.21$ to $1.05$---an almost exact variance match.
\end{itemize}

Aggregating over all seven scenarios with the symmetric fidelity error
$\lvert \log_2 J \rvert$, TFL reduces the geometric-mean fold error
from $7.6\times$ to $2.9\times$ and the median fold error from
$5.7\times$ ($2^{2.52}$) to $1.6\times$ ($2^{0.70}$).
\Cref{fig:jitterratio} visualises the per-test ratios on a log axis;
\cref{fig:tflcomparison} overlays the three distributions for
representative scenarios.

\begin{figure}[tb]
  \centering
  \includegraphics[width=\textwidth]{fig63_jitter_ratio.pdf}
  \caption[Jitter ratio, baseline vs.\ TFL]{Jitter ratio $J$ per
    scenario (log scale; $1.0$ = perfect variance fidelity). TFL moves
    every timer-driven scenario and T2 markedly toward $1$.}
  \label{fig:jitterratio}
\end{figure}

\begin{figure}[tb]
  \centering
  \includegraphics[width=\textwidth]{fig66_tfl_comparison.pdf}
  \caption[Distribution comparison hardware / baseline /
    TFL]{Interval distributions for hardware, baseline Renode, and
    TFL-enhanced Renode. TFL reshapes the emulated distributions
    toward the hardware envelope.}
  \label{fig:tflcomparison}
\end{figure}

\paragraph{Where TFL falls short.} Two scenarios regress. In T1 the
injected jitter is small relative to the (already suppressed) tick
distribution, leaving $J = 0.10$. In T3 the period-proportional jitter
model overshoots: hardware tick-rounding jitter is \emph{absolute}
(${\approx}\SI{400}{\micro\second}$ regardless of period), while the
TFL sleep-path perturbation scales with the period, producing
$\sigma = \SI{4.0}{\milli\second}$ at a \SI{200}{\milli\second} period
($J = 10$). The correct model for \srcS{4} sleep deviation is an
absolute, uniform $0$--$1$ tick term, not a proportional Gaussian---a
straightforward calibration fix identified by, and only visible
through, this hardware-grounded evaluation
(cf.\ \cref{sec:future:calibration}).

\section{Discussion and Threats to Validity}
\label{sec:eval:discussion}

\paragraph{Answer to RQ1.} Stock Renode's temporal deviation from
hardware is large, structured, and bidirectional: variance suppression
of $5\times$ for tick-quantised sleeps and variance inflation up to
$23\times$ for timer-driven events, with accurate means throughout.
Mean-based validation is therefore insufficient; distribution-level
metrics such as the jitter ratio are necessary.

\paragraph{Answer to RQ2.} A runtime-loaded peripheral layer can
recover a large share of temporal fidelity (geometric-mean fold error
$7.6\times \to 2.9\times$) without touching the emulator core or the
firmware logic. Residual error is dominated by the scheduling quantum
(\srcS{2}/\srcS{4}) and by calibration of the injected models.

\paragraph{Threats to validity.}
\begin{description}
  \item[Internal.] GPIO instrumentation adds \SI{261}{\nano\second}
        per marker (\cref{tab:instrumentation}), and the capture chain
        contributes \SI{12}{\nano\second} noise
        (\cref{tab:validation})---both negligible against measured
        effects. The TFL PRNG is seeded; all statistics are computed
        over full capture sets.
  \item[Construct.] The jitter ratio captures second-moment fidelity
        only. Mann--Whitney tests and the CDF/histogram comparisons
        (\cref{fig:perioddist,fig:cdf,fig:tflcomparison}) guard against
        distributions that match $\sigma$ but differ in shape.
  \item[External.] Results are demonstrated for one MCU family
        (STM32F1/Cortex-M3), one RTOS (Zephyr), and one emulator
        (Renode~1.15.3). The deviation taxonomy and the TFL mechanism
        are platform-generic, but the quantitative constants are not;
        transferring them requires re-calibration
        (\cref{sec:future:generalisation}).
\end{description}
